\chapter{Calculators}
\label{ch:calculators}

All calculators in \jns{} implement the ASE \code{Calculator} protocol
\cite{larsen2017ase}.  They can be used stand-alone --- geometry
optimisation with ASE, single-point scans, custom workflows --- without the
SCINE engine, and they are the building blocks the engine integration
(Chapter~7) resolves.

\begin{table}[htbp]
  \centering
  \small
  \begin{tabularx}{\linewidth}{@{}l l l X@{}}
    \toprule
    Class & Level & Backend & Role \\
    \midrule
    \code{SafeNNPCalculator} & L0 & NNP committee + ORCA fallback & exploration with run-time trust control \\
    \code{XTBCalculator} & L1 & GFN2-xTB \cite{bannwarth2019gfn2} & fast pre-screening, initial geometries \\
    \code{ORCACalculator} & L2 & ORCA r2SCAN-D4/def2-SVP & geometry optimisation and frequencies (production) \\
    \code{ORCACalculator} & L3 & ORCA wB97M-V/def2-TZVP & high-level single points \\
    \code{NNPCommittee} & --- & MACE-MP-0, MACE-MPA-0, AIMNet2-b973c & fast potential with a disagreement measure \\
    \bottomrule
  \end{tabularx}
  \caption{The calculators and their levels.  \code{CalculatorFactory}
  returns one per level; \code{NNPCommittee} is not a calculator itself but
  the predictor behind L0.}
\end{table}

\section{Conventions shared by all calculators}

\paragraph{Units.}
ASE conventions throughout: energy in eV, forces in eV/\AA.  The external
programs report Hartree and Hartree/bohr; the conversion is
$27.211386245$ for the energy and $F\,[\text{eV/\AA}] = -\nabla E \cdot
27.211386245 / 0.5291772109$ for the forces.

\paragraph{Properties.}
Every calculator in this version declares exactly \code{energy} and
\code{forces}.  There is no stress tensor, no Hessian and no charge output
from these classes.

\paragraph{Caching.}
Results follow the ASE protocol: one backend evaluation produces energy and
forces together, and repeated property requests on an unchanged \code{Atoms}
object reuse the cached values without a new evaluation.  To force a
re-evaluation, change the geometry or use a fresh calculator instance.  This
also means the counters of \code{SafeNNPCalculator} count evaluations, not
property accesses --- a distinction that matters when measuring cost or
fallback rates.

\paragraph{Scratch files and failures.}
External programs run in a temporary directory that is removed after the
call.  Passing \code{keep_files=True} together with \code{workdir} keeps
the input, output and gradient files there for inspection.  A backend that
exits non-zero or writes no result file raises \code{RuntimeError} naming
the output file.  Constructing a calculator never launches a program; the
backend runs on the first calculation.

\section{\code{XTBCalculator} (L1)}

GFN2-xTB \cite{bannwarth2019gfn2} through the \code{xtb} command line, for
fast pre-screening and initial geometries.

\begin{table}[htbp]
  \centering
  \small
  \begin{tabularx}{\linewidth}{@{}l l X@{}}
    \toprule
    Parameter & Default & Meaning \\
    \midrule
    \code{method} & \file{"GFN2-xTB"} & also \file{"GFN1-xTB"}, \file{"GFN0-xTB"}, \file{"GFNFF"} \\
    \code{charge} & \code{0} & total charge passed as \code{--chrg} \\
    \code{mult} & \code{1} & spin multiplicity; converted to \code{--uhf mult-1}, i.e.\ the number of unpaired electrons \\
    \code{keep_files} & \code{False} & keep the scratch files in \code{workdir} \\
    \code{workdir} & \code{None} & scratch directory (used when \code{keep_files} is true) \\
    \code{xtb_bin} & \code{XTB_BIN} or \file{xtb} & executable to call \\
    \bottomrule
  \end{tabularx}
  \caption{\code{XTBCalculator} arguments.}
\end{table}

\begin{codeblock}
from ase.build import molecule
from jnuscout.calculators import XTBCalculator

atoms = molecule("H2O")
atoms.calc = XTBCalculator(method="GFN2-xTB")
energy_ev = atoms.get_potential_energy()
forces = atoms.get_forces()          # shape (3, 3), eV/A
\end{codeblock}

Notes: the wrapper writes \file{mol.xyz}, runs xtb with
\code{--gfn}/\code{--chrg}/\code{--uhf}/\code{--grad}, reads the total energy
from the output and the gradient from the gradient file (Hartree/bohr,
converted to eV/\AA).  The executable must be installed
(Chapter~\ref{ch:installation}); a run that fails or writes no gradient
raises \code{RuntimeError}.

\section{\code{ORCACalculator} (L2/L3)}

Single-point energy and gradient from ORCA \cite{neese2022orca}.  L2 is the
production level for geometry optimisation and frequencies; L3 provides
high-level single points.  The keyword lines are defined in
\code{ORCA_KEYWORDS}:

\begin{table}[htbp]
  \centering
  \small
  \begin{tabularx}{\linewidth}{@{}l X@{}}
    \toprule
    Level & ORCA keywords \\
    \midrule
    \code{"L2"} & \code{r2SCAN D4 def2-SVP RIJCOSX def2/J TightSCF} \\
    \code{"L3"} & \code{wB97M-V def2-TZVP RIJCOSX AutoAux defgrid3 TightSCF} \\
    \bottomrule
  \end{tabularx}
  \caption{Method and basis per level (r2SCAN \cite{furness2020r2scan},
  D4 \cite{caldeweyher2019d4}, wB97M-V \cite{mardirossian2016wb97mv}).}
\end{table}

\begin{table}[htbp]
  \centering
  \small
  \begin{tabularx}{\linewidth}{@{}l l X@{}}
    \toprule
    Parameter & Default & Meaning \\
    \midrule
    \code{level} & \file{"L2"} & \file{"L2"} or \file{"L3"}; anything else raises \code{ValueError} \\
    \code{nprocs} & \code{8} & the \%\code{pal nprocs} directive in the input \\
    \code{maxcore} & \code{4000} & the \%\code{maxcore} directive, memory in MB per core \\
    \code{keep_files} & \code{False} & keep \file{.inp}/\file{.out}/\file{.engrad} in \code{workdir} \\
    \code{workdir} & \code{None} & scratch directory (used when \code{keep_files} is true) \\
    \code{orca_bin} & \code{ORCA_BIN} or \file{orca} & executable to call \\
    \bottomrule
  \end{tabularx}
  \caption{\code{ORCACalculator} arguments.}
\end{table}

\begin{codeblock}
from ase import Atoms
from jnuscout.calculators import ORCACalculator

atoms = Atoms("H2O", positions=[[0.0, 0.0, 0.0], [0.0, 0.0, 0.96], [0.0, 0.93, -0.24]])
atoms.calc = ORCACalculator(level="L2", nprocs=4, maxcore=2000)
energy_ev = atoms.get_potential_energy()
forces = atoms.get_forces()          # eV/A
\end{codeblock}

Notes:
\begin{itemize}
  \item The charge and multiplicity of the ORCA input are taken from the
  \code{initial_charges} and \code{initial_magnetic_moments} arrays of the
  \code{Atoms} object (their sums); with neither array present the input is
  neutral and singlet.
  \item L2 uses RIJCOSX with the def2/J auxiliary basis.  r2SCAN is a pure
  meta-GGA without Hartree--Fock exchange, so ORCA automatically downgrades
  RIJCOSX to Split-RI-J and prints a warning; the warning is harmless and
  there is no separate \code{RI-J} keyword.
  \item L3 uses an AutoAux-generated auxiliary basis.  The originally
  specified def2/JK set was not available on the deployment host (no local
  auxiliary-basis file and no network access for the automatic download), so
  AutoAux is the tested default.
  \item The input asks for \code{Engrad}; the calculator parses the
  \file{.engrad} file for energy and gradient and verifies the atom count
  against the structure.
  \item ORCA is launched in an adjusted environment; if the binary cannot
  find its libraries or \file{mpirun}, see the notes in
  Chapter~\ref{ch:installation}.
\end{itemize}

\section{\code{NNPCommittee}}

The committee is the prediction backend of L0: three foundation models from
different architecture families --- MACE-MP-0 and MACE-MPA-0
\cite{batatia2022mace,batatia2025mpa} and AIMNet2-b973c
\cite{anstine2025aimnet2} --- whose disagreement is the uncertainty signal.
Members from the same family (differently seeded fine-tunes) predict almost
the same and cannot serve as a measure, which is why the families are mixed.

\begin{table}[htbp]
  \centering
  \small
  \begin{tabularx}{\linewidth}{@{}l l X@{}}
    \toprule
    Member & Backend & Model identifier \\
    \midrule
    \code{MACE-MP-0} & \code{mace} & \code{MACE_MP0_PATH}, default \file{medium} \\
    \code{MACE-MPA-0} & \code{mace} & \code{MACE_MPA0_PATH}, default the cached medium model \\
    \code{AIMNet2-b973c} & \code{aimnet2} & \file{aimnet2_b973c} (weights resolved by the \file{aimnet2} package) \\
    \bottomrule
  \end{tabularx}
  \caption{Default committee members.  The AIMNet2 base is the \code{b973c}
  variant; the \code{wb97m_0} variant showed a force-directionality defect
  against ORCA references and is not used.}
\end{table}

\begin{table}[htbp]
  \centering
  \small
  \begin{tabularx}{\linewidth}{@{}l l X@{}}
    \toprule
    Argument & Default & Meaning \\
    \midrule
    \code{members} & the three members above & list of \code{(name, backend, model_id)} tuples \\
    \code{device} & \file{"cuda"} & \file{"cuda"} or \file{"cpu"}, passed to the backends \\
    \code{energy_baselines} & \code{None} & per-member per-atom energy offsets in eV, from calibration \\
    \code{energy_family} & \file{"mace"} & backend family used for the energy disagreement; \code{None} picks the largest family \\
    \bottomrule
  \end{tabularx}
  \caption{\code{NNPCommittee} arguments.}
\end{table}

\begin{codeblock}
from ase.build import molecule
from jnuscout.calculators import NNPCommittee

committee = NNPCommittee(device="cuda")
report = committee.compute(molecule("H2O"))
print(report["uq_force"], report["uq_energy"])   # eV/A and eV/atom
print(report["energy_mean"])                     # eV (ensemble)
print(report["per_member_energy"])               # per-member raw energies
\end{codeblock}

\begin{readbox}[Disagreement definitions]
For a committee of $M$ members on a system of $N$ atoms, with $F_{i,j}$ the
force on atom $j$ from member $i$:
\[
uq_F = \max_{j} \max_{i} \lVert F_{i,j} - \bar{F}_{j} \rVert
\quad[\text{eV/\AA}],
\qquad
uq_E = \frac{\operatorname{std}_i(E_i)}{N}
\quad[\text{eV/atom}].
\]
The force disagreement is a maximum over atoms, not a mean: in reacting
geometries the failure is usually one atom with a wrong force direction,
which averaging would hide.  The energy disagreement is computed only among
members of one backend family (\code{energy_family}), because absolute
energies of different families have different zeros --- MACE reports
atomization energies, AIMNet2 absolute total energies --- so subtracting
across families is not meaningful.
\end{readbox}

The report returned by \code{disagreement()} and \code{compute()} contains
\code{uq_energy}, \code{uq_force}, \code{energy_mean} (eV, ensemble),
\code{forces_mean} (\AA$\times$3 array, eV/\AA), \code{energy_spread} and
\code{energy_spread_raw} (after and before baseline alignment),
\code{baselines_aligned} (bool), \code{per_member_energy}, and --- from
\code{compute()} --- \code{per_member_forces}.  \code{compute()} also
stores the report in \code{last_report}.

Notes:
\begin{itemize}
  \item Member calculators are built lazily on the first prediction and
  cached per member; the first call therefore includes the model load.
  \item Energy baselines are per-atom offsets.  Without calibration,
  \code{uq_energy} is a diagnostic on raw absolute energies and has no
  physical meaning.  \code{calibrate_baselines}\allowbreak\code{(atoms_list)} estimates the
  offsets from the per-model mean energies over a common set of
  configurations; the alignment is approximate, and any configuration-dependent
  bias of a member leaks into \code{uq_energy}.  The force disagreement is
  unaffected by constant offsets, which is why it is the default criterion.
  \item Degenerate inputs are reported as \code{logging} warnings, not
  errors: prediction keys that do not match the members, an
  \code{energy_family} with no available members, and members missing from
  \code{energy_baselines}.  Fewer than two predictions raise
  \code{ValueError}.
\end{itemize}

\section{\code{SafeNNPCalculator} (L0)}

The committee plus a calibrated threshold: below it the ensemble mean is
returned, at or above it the structure is re-evaluated with ORCA L2.  This
is the calculator the engine uses during exploration, and the mechanism the
package is built around.  The fallback idea is related to, but not the same
as, the smoothly added xTB guardrail of Staub et al.\ \cite{staub2023}:
there it participates in every prediction with no threshold, whereas here
the decision is explicit and calibrated.

\begin{table}[htbp]
  \centering
  \small
  \begin{tabularx}{\linewidth}{@{}l l X@{}}
    \toprule
    Argument & Default & Meaning \\
    \midrule
    \code{theta_force} & (required) & force-disagreement threshold in eV/\AA; from calibration \\
    \code{theta_energy} & \code{None} & energy threshold in eV/atom; \code{None} = force criterion only \\
    \code{fallback_level} & \file{"L2"} & level passed to the fallback \code{ORCACalculator} \\
    \code{committee} & fresh \code{NNPCommittee} & inject a prepared committee \\
    \code{device} & \file{"cuda"} & passed to the default committee \\
    \code{use_or_rule} & \code{True} & trigger on either criterion (see below) \\
    \code{fallback_kwargs} & \code{None} & extra keywords for the fallback calculator, e.g.\ \code{{"nprocs": 8}} \\
    \bottomrule
  \end{tabularx}
  \caption{\code{SafeNNPCalculator} arguments.}
\end{table}

The decision rule, with $\theta_F$ = \code{theta_force} and $\theta_E$ =
\code{theta_energy}:

\begin{table}[htbp]
  \centering
  \small
  \begin{tabularx}{\linewidth}{@{}l X@{}}
    \toprule
    Configuration & Fallback when \\
    \midrule
    \code{theta_energy} is \code{None} & $uq_F \ge \theta_F$ \\
    \code{use_or_rule=True} (default) & $uq_F \ge \theta_F$ or $uq_E \ge \theta_E$ (conservative) \\
    \code{use_or_rule=False} & both thresholds are reached (AND rule) \\
    \bottomrule
  \end{tabularx}
  \caption{Fallback conditions.  The comparison is \code{>=}; a disagreement
  exactly equal to the threshold triggers the fallback.}
\end{table}

\begin{codeblock}
from ase.io import read
from jnuscout.calculators import SafeNNPCalculator

calc = SafeNNPCalculator(theta_force=0.30, theta_energy=0.05,
                         device="cuda", fallback_kwargs={"nprocs": 8})
atoms = read("molecule.xyz")
atoms.calc = calc
energy_ev = atoms.get_potential_energy()

print(calc.n_fallback, "/", calc.n_calls)   # fallbacks / evaluations
print(calc.stats())                         # dict with the fallback rate
print(calc.last_path)                       # "nnp" or "fallback"
print(calc.last_uq)                         # full report of the last evaluation
\end{codeblock}

Notes:
\begin{itemize}
  \item \code{theta_force} has no default.  It must come from the
  calibration protocol (path~1 of Chapter~\ref{ch:getting-started}); the
  factory's L0 requires it as well.
  \item The fallback calculator is constructed on the first fallback
  (using \code{fallback_level} and \code{fallback_kwargs}) and reused
  afterwards, so ORCA is not started unless a fallback occurs.
  \item The committee runs once per calculation.  On the NNP path the
  returned energy and forces are the ensemble mean; on the fallback path the
  ORCA values replace them entirely, and the NNP result for that geometry is
  discarded.
  \item \code{n_calls} counts committee evaluations, \code{n_fallback} the
  subset routed to ORCA; both follow the ASE caching semantics described
  above.
\end{itemize}

\section{\code{CalculatorFactory}}

One call per level, returning the calculator described in the sections
above:

\begin{table}[htbp]
  \centering
  \small
  \begin{tabularx}{\linewidth}{@{}l l X@{}}
    \toprule
    Level & Returns & Notes \\
    \midrule
    \code{"L0"} & \code{SafeNNPCalculator(**kwargs)} & \code{theta_force} is required; no default threshold \\
    \code{"L1"} & \code{XTBCalculator(...)} & consumes \code{method} (default \file{"GFN2-xTB"}) \\
    \code{"L2"} & \code{ORCACalculator(level="L2", ...)} & \\
    \code{"L3"} & \code{ORCACalculator(level="L3", ...)} & \\
    \bottomrule
  \end{tabularx}
  \caption{\code{CalculatorFactory.create(level, **kwargs)} routing.}
\end{table}

\begin{codeblock}
from jnuscout.calculators import CalculatorFactory

safe = CalculatorFactory.create("L0", theta_force=0.30)
cheap = CalculatorFactory.create("L1", charge=-1, mult=2)
prod = CalculatorFactory.create("L2", nprocs=8, maxcore=4000)
\end{codeblock}

Notes: \code{create} is a static method, callable on the class or an
instance.  By design L0 has no default threshold, so
\code{CalculatorFactory.create("L0")} raises \code{TypeError}; an unknown
level raises \code{ValueError}.  Construction itself never launches a
backend.
